\section{实验设计}
\subsection{数据、参数依据与分阶段预算}
主实验只在500节点欧氏TSP上训练，1000节点实例仅用于冻结后的规模迁移测试。实例、参考路径和面板归属由现有数据清单固定，清单将该合成实例族记为uniform。现存元数据没有完整记录生成器及父实例关系，本文不额外假定未记录的生成过程。参考路径是数据集提供的比较基准，未逐例独立认证最优性。对路径$T$，评价指标为
\begin{equation}
 \gap(T)=100\left(\frac{C(T)}{C_{\mathrm{ref}}}-1\right)\%,
\end{equation}
其中$C$按连续欧氏距离计算。文中称其为\emph{参考差距}，而非已认证最优差距。标签仅由外部评价器读取，以高精度求和计算长度；在线求解器不读取标签或标签路径。

蚂蚁数采用2022 FACO实验的规模规则$m=64\lceil\sqrt n/16\rceil$\cite{skinderowicz2022}，TSP500和TSP1000均为128。主要／备用候选数16／64、局部搜索候选数20、$\beta=1$、信息素保留系数0.5、$p_{\mathrm{best}}=0.1$、阶段最优使用概率0.01在GPU方法间固定。MNE8作为FACO式默认强度对照，MNE16用于排除最大强度本身带来的收益。GP种群、节点上限和算子概率是控制搜索规模的实验设计选择，并未经过全面超参数优化。

训练迭代数允许少于测试，但不按时间截断。先前开发标定比较24个固定或规则控制器，使用32个开发实例与5个ACO种子，在候选$H\in\{100,250,500,1000,2500,5000\}$上要求保留至少95\%的相对初始化改善，且与$H=5000$的控制器排名Spearman相关至少0.9。TSP500中，$H=1000$是首个同时达标的候选，改善保留率为98.73\%，排名相关为0.9383。此依据来自开发控制器，并不保证进化GP在短、长预算下排名一致。

表\ref{tab:protocol}给出正式轮次预算。三种GP种子为1103、2207、3313，每种表示各运行一次，共6次。所有运行在同一代共享实例与ACO随机种子面板，训练面板随代变化，50个面板合计覆盖794个不同训练实例。监控、选模和测试均不含这些训练实例，快速验证与完整验证的实例也不重叠。每个个体每代$16\times1\times128\times1000=2{,}048{,}000$ FE，每次50代训练为$13{,}107{,}200{,}000$ FE；六次合计$78{,}643{,}200{,}000$ FE，不含监控、选模和测试。

\begin{table}[tb]
\caption{正式实验流程。每次求解均使用128只蚂蚁；快速与完整验证使用不同面板。}
\label{tab:protocol}
\small
\begin{tabularx}{\textwidth}{@{}lrrY@{}}
\toprule
阶段 & 实例$\times$ACO种子 & $H$ & 评价对象与用途\\ \midrule
训练 & $16\times1$ & 1000 & 每代128个体，50代，面板随代变化\\
开发监控 & $16\times1$ & 1000 & 第5、10、$\ldots$、50代冠军；固定面板\\
快速验证 & $32\times1$ & 1000 & 历代冠军与末代种群，去重后最多177个\\
完整验证 & $64\times3$ & 5000 & 每次运行快速验证前4名，选出1个\\
TSP500测试 & $128\times10$ & 5000 & 六个冻结控制器及四种对照\\
TSP1000测试 & $128\times10$ & 5000 & 原样迁移，不重新进化或选模\\
\bottomrule
\end{tabularx}
\end{table}

监控和快速验证的ACO种子为17；完整验证为17、29、41；测试为17、29、41、53、67、79、97、109、127、149。划分以实例清单为准，重复使用种子编号不意味着重复使用测试实例。六个控制器全部冻结后才执行正式测试。本次整理没有新增ACO评价，也没有根据测试结果修改方法。

\subsection{对照与统计方法}
四种对照分别为：\FACO{}，保留24条初始化路径的原生2022 FACO；\Ueight{}，GPU底座、MNE8、全体节点均匀起点且不重启；\Usixteen{}，仅将前者强度改为16；\Rsixteen{}，GPU底座、MNE16、每迭代随机生成的区域0起点且不重启。区域0对照可以检验局部起点限制是否已足以解释收益。GPU静态对照与GP使用相同构造、局部搜索及数值后端。

\FACO{}在合成数据上使用连续FP64距离，并将依赖距离取整的KD-tree邻域实现替换为近邻列表；对恒等2-opt移动赋零增益，并排除边集合不变的3-opt移动，以避免浮点误差触发虚假改善循环。它应理解为原生FACO的连续距离适配版。另有保留原始整数距离的TSPLIB\cite{reinelt1991}复现实验，详见补充材料；二者不合并为同一结果。

质量均值先对ACO种子求平均，再对实例和GP种子求平均。主检验以128个测试实例为配对单位，每个实例先平均10个ACO种子与3个GP种子，采用双侧Wilcoxon符号秩检验\cite{wilcoxon1945}。TSP500预定的9项比较为两种GP各对四种对照及GP之间比较，统一进行Holm校正\cite{holm1979}。95\%区间采用10\,000次分层配对bootstrap，重采样GP种子、实例及实例内ACO种子，统计种子为91001。定义改善$\Delta=\gap_{\text{对照}}-\gap_{\mathrm{GP}}$，以百分点报告，正值表示GP更好。TSP1000区间作为迁移效果的描述性证据，不额外扩展主检验族。3个GP种子仍限制了对进化随机性的估计精度。
